\documentclass[11pt,a4paper]{article}

% Language and fonts (compile with XeLaTeX)
\usepackage{fontspec}
\usepackage[english]{babel}

\usepackage[margin=2.5cm]{geometry}
\usepackage{amsmath, amssymb}
\usepackage{graphicx}
\usepackage{enumitem}
\usepackage[hidelinks]{hyperref}

% Assignment details
\newcommand{\course}{Course name}
\newcommand{\assignment}{Assignment 1}
\newcommand{\student}{First Last}
\newcommand{\instructor}{Instructor name}

% Numbered question and answer environment
\newcounter{question}
\newcommand{\question}[1]{%
  \stepcounter{question}%
  \section*{Question \thequestion\ifx\relax#1\relax\else: #1\fi}}
\newenvironment{answer}{\par\noindent\textbf{Answer.}\ }{\hfill$\blacksquare$\par}

\begin{document}

\noindent
\begin{minipage}[t]{0.6\textwidth}
  \textbf{\course}\\
  \assignment
\end{minipage}%
\begin{minipage}[t]{0.4\textwidth}
  \raggedleft
  \student\\
  Instructor: \instructor\\
  \today
\end{minipage}
\par\medskip\hrule\medskip

\question{Derivatives}
Compute $\dfrac{d}{dx}\left(x^2 \sin x\right)$.

\begin{answer}
  By the product rule,
  \[
    \frac{d}{dx}\left(x^2 \sin x\right) = 2x \sin x + x^2 \cos x.
  \]
\end{answer}

\question{}
State two properties of the definite integral.

\begin{answer}
  \begin{enumerate}[label=(\alph*)]
    \item Linearity: $\int_a^b (\alpha f + \beta g) = \alpha \int_a^b f + \beta \int_a^b g$.
    \item Additivity: $\int_a^b f = \int_a^c f + \int_c^b f$.
  \end{enumerate}
\end{answer}

\end{document}
